\documentclass[10pt,a4paper]{amsart}
\usepackage[margin=19mm]{geometry}
\usepackage{amsmath,amssymb,mathtools}
\usepackage[scr=rsfs,cal=euler]{mathalfa}
\usepackage{xfrac}
\usepackage[unicode,hidelinks,bookmarks=false]{hyperref}
\newcommand{\ie}{i.e.\ }
\newcommand{\cf}{cf.\ }
\newcommand{\resp}{respectively\ }
\newcommand{\Subsection}{\subsection}
\providecommand{\nobreakdash}{\nobreak\hbox{-}}
\setlength{\emergencystretch}{2em}
\multlinegap0pt
\usepackage{bm}
\usepackage{bbm}
\usepackage{dsfont}
\usepackage{xspace}
\usepackage{cancel}
\usepackage{mathtools}
\usepackage{amsthm}
\makeatletter
\@ifundefined{lemma}{\newtheorem{lemma}{Lemma}}{}
\@ifundefined{theorem}{\newtheorem{theorem}{Theorem}}{}
\makeatother
\title[The Number of Parts in the (Distinct) Partitions With Parts ]{The Number of Parts in the (Distinct) Partitions With Parts From a Set}
\author{A. David Christopher}
\date{Source: arXiv:2509.22618v1 (CC BY 4.0)}
\begin{document}
\maketitle

\noindent\emph{Source and licence.} The Number of Parts in the (Distinct) Partitions With Parts From a Set. Version arXiv:2509.22618v1 (CC BY 4.0), distributed under the Creative Commons Attribution 4.0 International licence, as recorded in the arXiv licence metadata for this exact version and shown on the abstract page https://arxiv.org/abs/2509.22618v1. Full rights notice: licenses/arxiv-2509.22618-CC-BY-4.0.txt.\par\smallskip

\noindent\textit{Editorial scope.} Counted unit: the Theorem whose source label is \texttt{mest} in the article (source0.tex lines 550--555, source label \texttt{mest}), delivered together with the article's own complete original proof of it (source0.tex lines 556--703, one \texttt{proof} environment of 148 lines). No mathematics is written, repaired or re-derived: statement and proof are the article's own text line for line. Required context retained uncounted: roverA (lines 446--452); lemma2 (lines 546--548). Every definition, display and label of those lines is retained unchanged. Omitted from this package and not redistributed: 1--445, 453--545, 549--549, 704--739. Nothing inside a retained excerpt is omitted or altered: every retained body is delivered exactly as the article's lines are written, so no omission receipt of any kind accompanies this package. Citations: the retained excerpts carry no citation command at all, so no bibliography entry of the article is transcribed and no reference is redistributed. Reference adaptation: none was needed and none was made; every reference inside the delivered unit resolves inside it, so no unresolved reference or citation is produced. Printed slips kept literal: no printed word, symbol, hypothesis or number of the counted statement or of its proof was altered. Layout and class: the article's own class is \texttt{article} with packages including amsmath, amssymb, amsthm, hyperref, geometry; the wrapper is the lane's pinned \texttt{amsart} editorial wrapper at 10pt on A4, which restates the theorem-like environments the retained text needs and adds the article's own macro definitions that the retained text needs (none), with extra wrapper packages (bm, bbm, dsfont, xspace, cancel, mathtools). The delivered numbering is the unit's own. Assets: no retained span contains a figure environment, a graphics-inclusion command, a PDF-page inclusion, a table or a TikZ picture; no image file is copied into this package, the package declares no asset dependency and no image file is redistributed with it. Licence: version arXiv:2509.22618v1 of this article is distributed under the Creative Commons Attribution 4.0 International licence, as recorded in the arXiv licence metadata for this exact version and shown on the abstract page https://arxiv.org/abs/2509.22618v1. A full rights notice is delivered at licenses/arxiv-2509.22618-CC-BY-4.0.txt. Characters: every retained excerpt is pure ASCII, so the delivered engine, XeLaTeX with its default Latin Modern text font, reports no missing character.
\begin{theorem}\label{roverA}
Let $n$ be a positive integer and let $A$ be a set of positive integers. For $s\in A$, we have
\begin{equation}
N_A^p(n)-N_A^p(n-s)=p_A(n-s)+N_{A\setminus\{s\}}^p(n)
\end{equation}
with $p_A(0)=1$.
\end{theorem}

\begin{lemma}\label{lemma2}
Let $l$ be a non-negative integer and let $A=\{a_1,a_2,\cdots ,a_k\}$ be a set of positive integers such that $\gcd(A)=1$. For each $r\in\{0,1,\cdots,a_1a_2\cdots a_k-1\}$, the term $p_A(a_1a_2\cdots a_kl+r)$ is a polynomial in $l$ of degree $k-1$ with the leading coefficient $\frac{(a_1a_2\cdots a_k)^{k-2}}{(k-1)!}$.
\end{lemma}

\begin{theorem}
Let $A=\{a_1,a_2,\cdots ,a_k\}$ be a set of positive integers such that $\gcd(A)=1$. Then we have
\begin{equation}\label{mest}
N_A^p(n)\sim\frac{1}{k!}\frac{\frac{1}{a_1}+\frac{1}{a_2}+\cdots +\frac{1}{a_k}}{a_1a_2\cdots a_k}n^k.
\end{equation}
\end{theorem}

\begin{proof}
The crux of this proof is to show that $N_A^p(a_1a_2\cdots a_kl+r)$ is a polynomial in $l$ of degree $k$ with the leading coefficient \[\frac{\sum_{i=1}^{k}\left(a_i^{k-2}\prod_{\substack{j\neq i\\ 0\leq j\leq k}}a_j^{k-1}\right)}{k!}\] for each $r\in\{0,1,\cdots ,a_1a_2\cdots a_k-1\}$. Once this is established, then we will have
\[\lim_{l\to\infty}\frac{N^p_A(a_1a_2\cdots a_kl+r)}{(a_1a_2\cdots a_kl+r)^{k}}=\frac{1}{k!}\frac{\frac{1}{a_1}+\frac{1}{a_2}+\cdots +\frac{1}{a_k}}{a_1a_2\cdots a_k}
\]
for each $r\in\{0,1,\cdots ,a_1a_2\cdots a_k-1\}$, which is equivalent to the asymptotic estimate (\ref{mest}).

We will establish this main claim using induction over $|A|=k$. Since the term $a_i^{k-2}$ is involved in the leading coefficient, we take $k=3$ as the initial case for induction. Nevertheless the cases $k=1,2$ are needed to realise the case $k=3$. 

When $|A|=1$ the only way out is $a_1=1$. In this case we have $N_A^p(n)=n$, and the targeted estimate follows. To proceed further, we need the following observation: 
\[
N_{\{a\}}^p(n)=\begin{cases}\frac{n}{a} &\text{ if }a\mid n;\\
0&\text{ otherwise. }
\end{cases}
\]
Assume  $|A|=2$. That is, $A=\{a_1,a_2\}$ with $\gcd(a_1,a_2)=1$.
Fix $r\in\{1,2,\cdots ,a_1\}$. Applying Theorem \ref{roverA} $a_1$ times, we get
\begin{align*}
N_A^p(a_1a_2l+r)-N_A^p(a_1a_2(l-1)+r)&=\sum_{m=1}^{a_1}p_A(a_1a_2l+r-ma_2)\\ &\ \ \ +\sum_{t=0}^{a_1-1}N_{\{a_1\}}^p(a_1a_2l+r-ta_2).
\end{align*}
Since $\gcd(a_1,a_2)=1$, the congruence equation $a_2t\equiv r\pmod {a_1}$ has a unique solution modulo $a_1$, say $t^*$. This gives
\begin{align*}
\sum_{t=0}^{a_1-1}N_{\{a_1\}}^p(a_1a_2l+r-ta_2)&=\frac{a_1a_2l+r-t^*a_2}{a_1}\\ 
&=a_2l+\frac{r-t^*a_2}{a_1}\\
&=a_2l+k_1
\end{align*}
for each $r\in\{0,1,\cdots ,a_1-1\}$, where $k_1$ is an integer constant.

Next, we analyze the sum $\sum_{m=1}^{a_1}p_A(a_1a_2l+r-ma_2)$. Let $m\in\{1,2,\cdots ,a_1\}$. \newline 
Case i. Assume $r-ma_2<0$. In this case, consider the term $p_A(a_1a_2l+r-ma_2)$, which can be written as $p_A(a_1a_2(l-1)+a_1a_2+r-ma_2)$. Here we note that $0\leq a_1a_2+r-ma_2\leq a_1a_2-1$. Now, in view of Lemma \ref{lemma2}, we have
\begin{align*}
p_A(a_1a_2(l-1)+a_1a_2+r-ma_2)=(l-1)+c_m,
\end{align*}
where $c_m$ is an integer constant.\newline 
Case ii. Assume $r-ma_2\geq 0$. Then by Lemma \ref{lemma2}, we have $p_A(a_1a_2l+r-ma_2)=l+d_m$,
where $d_m$ is an integer constant.

In both the cases, $p_A(a_1a_2l+r-ma_2)=l+b_m$ for some constant $b_m$. This gives
\begin{align*}
\sum_{m=1}^{a_1}p_A(a_1a_2l+r-ma_2)&=\sum_{m=1}^{a_1}(l+b_m)\\
&=a_1l+k_2,
\end{align*}
where $k_2=\sum_{m=1}^{a_1}b_m$ is an integer constant. \newline 
Consequently,
\[N_A^p(a_1a_2l+r)-N_A^p(a_1a_2(l-1)+r)=(a_1+a_2)l+c,
\]
where $c=k_1+k_2$ is an integer constant.

This gives
\[N_A^p(a_1a_2l+r)=(a_1+a_2)\left(\frac{l^2+l}{2}\right)+cl+N_A^p(r).
\]
Subsequently, we have
\[\lim_{l\to\infty}\frac{N_A^p(a_1a_2l+r)}{(a_1a_2l+r)^2}
=\frac{\frac{1}{a_1}+\frac{1}{a_2}}{2!(a_1a_2)}.
\]
Since this is true for each $r\in\{0,1,\cdots ,a_1a_2-1\}$, the targeted estimate follows for the case $k=2$.

Now we have the initial case of the induction, that is, $k=3$. Assume $A=\{a_1,a_2,a_3\}$ with $\gcd(a_1,a_2,a_3)=1$. Fix $r\in\{0,1,\cdots ,a_1a_2a_3-1\}$. 

Let $d=\gcd(a_1,a_2)$. Applying Theorem \ref{roverA} $a_1a_2$ times, we get
\begin{align*} 
&N_A^p(a_1a_2a_3l+r)-N_A^p(a_1a_2a_3(l-1)+r)\\ 
&=\sum_{t=1}^{a_1a_2}p_A(a_1a_2a_3l+r-ta_3)+\sum_{m=0}^{a_1a_2-1}N_{A\setminus \{a_3\}}^p(a_1a_2a_3l+r-ma_3)\ \ \ \ \ \ \ \ \ \ \ \ \ \ \ \ \ \ \ \ \ \ \\ 
&=\sum_{t=1}^{a_1a_2}p_A(a_1a_2a_3l+r-ta_3)+\sum_{\substack{0\leq m\leq a_1a_2-1\\ d\mid r-ma_3}}N_{ \{a_1,a_2\}}^p\left(\frac{a_1a_2a_3l}{d}+\frac{r-ma_3}{d}\right)
\end{align*}
\begin{align}\label{eqn1}
=\sum_{t=1}^{a_1a_2}p_A(a_1a_2a_3l+r-ta_3)+\sum_{\substack{0\leq m\leq a_1a_2-1\\ d\mid r-ma_3}}N^p_{\{\frac{a_1}{d},\frac{a_2}{d}\}}\left(\frac{a_1}{d}\frac{a_2}{d}\frac{a_3d^2}{d}l+\frac{r-ma_3}{d}\right).\ \ 
\end{align}
If $r-ta_3<0$, then by Lemma \ref{lemma2}, $p_A(a_1a_2a_3(l-1)+a_1a_2a_3+r-ta_3)$ is a polynomial in $l-1$ of degree $2$ with the leading coefficient $\frac{a_1a_2a_3}{2!}$. If $r-ta_3\geq 0$, then again by Lemma \ref{lemma2}, $p_A(a_1a_2a_3l+r-ta_3)$ is a polynomial in $l$ of degree $2$ with the leading coefficient $\frac{a_1a_2a_3}{2!}$. We observe that, in both the cases, $p_A(a_1a_2a_3l+r-ta_3)$ is a polynomial in $l$ of degree $2$ with the leading coefficient $\frac{a_1a_2a_3}{2!}$. This in turn gives $\sum_{t=1}^{a_1a_2}p_A(a_1a_2a_3l+r-ta_3)$ is a polynomial in $l$ of degree 2 with the leading coefficient $\frac{a_1^2a_2^2a_3}{2!}$.  

Since $\gcd(d,a_3)=1$,  the congruence equation $a_3m\equiv r\pmod d$ has a unique solution modulo $d$. Consequently, there exists a unique integer   $m_i\in\{id+0,id+1,\cdots ,id+(d-1)\}$ such that $a_3m_i\equiv r\pmod d$ for each $0\leq i\leq\frac{a_1a_2}{d}-1$. Based on these observations, we can write

\[
\sum_{\substack{0\leq m\leq a_1a_2-1\\ d\mid r-ma_3}}N^p_{\{\frac{a_1}{d},\frac{a_2}{d}\}}\left(\frac{a_1}{d}\frac{a_2}{d}\frac{a_3d^2}{d}l+\frac{r-ma_3}{d}\right)\ \ \ \ \ \ \ \ \ \ \ \ \ \ \ \ \ \ \ \ \ \ \ \ \ \ \ \ 
\]


\begin{align*}
&=\sum_{i=0}^{\frac{a_1a_2}{d}-1}N^p_{\{\frac{a_1}{d},\frac{a_2}{d}\}}\left(\frac{a_1}{d}\frac{a_2}{d}\frac{a_3d^2}{d}l+\frac{r-m_ia_3}{d}\right)\\
&=\sum_{i=0}^{\frac{a_1a_2}{d}-1}N^p_{\{\frac{a_1}{d},\frac{a_2}{d}\}}\left(\frac{a_1}{d}\frac{a_2}{d}\left(a_3dl+q_i\right)+r_i\right),
\end{align*}
where $q_i$ and $r_i$ were uniquely determined (in view of division algorithm) satisfying the relation:$\frac{r-m_ia_3}{d}=\frac{a_1}{d}\frac{a_2}{d}q_i+r_i$. Since $\gcd(\frac{a_1}{d},\frac{a_2}{d})=1$, from the case $|A|=2$, we have that \[N^p_{\{\frac{a_1}{d},\frac{a_2}{d}\}}\left(\frac{a_1}{d}\frac{a_2}{d}\left(a_3dl+q_i\right)+r_i\right)\] is a polynomial in $a_3dl+q_i$ of degree 2 with the leading coefficient $\frac{\frac{a_1}{d}+\frac{a_2}{d}}{2!}$. This in turn gives that $N^p_{\{\frac{a_1}{d},\frac{a_2}{d}\}}\left(\frac{a_1}{d}\frac{a_2}{d}\left(a_3dl+q_i\right)+r_i\right)$ is a polynomial in $l$ of degree 2 with the leading coefficient $\frac{\frac{a_1}{d}+\frac{a_2}{d}}{2!}a_3^2d^2$. Consequently, \[\sum_{i=0}^{\frac{a_1a_2}{d}-1}N^p_{\{\frac{a_1}{d},\frac{a_2}{d}\}}\left(\frac{a_1}{d}\frac{a_2}{d}\left(a_3dl+q_i\right)+r_i\right)\] is a polynomial in $l$ of degree 2 with the leading coefficient 
$
\frac{a_1a_2}{d}\frac{\frac{a_1}{d}+\frac{a_2}{d}}{2!}a_3^2d^2
=\frac{a_1^2a_2a_3^2+a_1a_2^2a_3^2}{2!}$.

Now in view of Equation (\ref{eqn1}), we have that:
$
N_A^p(a_1a_2a_3l+r)-N_A^p(a_1a_2a_3(l-1)+r)
$ is a polynomial in $l$ of degree 2 with the leading coefficient $\frac{a_1^2a_2^2a_3+a_1^2a_2a_3^2+a_1a_2^2a_3^2}{2!}$.

From this we infer that $N_A^p(a_1a_2a_3l+r)$ is a polynomial in $l$ of degree 3. Now if one assumes $c_3$ to be the leading coefficient of $N_A^p(a_1a_2a_3l+r)$, then in accordance with the previous observations, we have
\[
3c_3=\frac{a_1^2a_2^2a_3+a_1^2a_2a_3^2+a_1a_2^2a_3^2}{2!}.
\]
This gives 
\[c_3=\frac{a_1^2a_2^2a_3+a_1^2a_2a_3^2+a_1a_2^2a_3^2}{3!}.
\]
Hence, the assertion is true for $k=3$. 

Assume that the assertion is true up to some $k-1\geq 3$. Now we shall prove the assertion for $k$. 

Let $A=\{a_1,a_2,\cdots ,a_k\}$ be a set of positive integers such that $\gcd(a_1,a_2,\cdots ,a_k)$\ $=1$. Let $d=\gcd(a_1,a_2,\cdots ,a_{k-1})$. Fix $r\in\{0,1,\cdots ,a_1a_2\cdots a_k-1\}$. Repeated application of Theorem \ref{roverA} $a_1a_2\cdots a_{k-1}$ times gives

$N_A^p(a_1a_2\cdots a_kl+r)-N_A^p(a_1a_2\cdots a_k(l-1)+r)$
\begin{align*}
&=\sum_{i=1}^{a_1a_2\cdots a_{k-1}}p_A(a_1a_2\cdots a_kl+r-ia_k)\\ &\ \ \ \ \ +\sum_{\substack{0\leq m\leq a_1a_2\cdots a_{k-1}-1\\ d|r-ma_k}}N_{A\setminus\{a_k\}}^p(a_1a_2\cdots a_kl+r-ma_k)\\
&=\sum_{i=1}^{a_1a_2\cdots a_{k-1}}p_A(a_1a_2\cdots a_kl+r-ia_k)\\ &\ \ \ \ \ +\sum_{\substack{0\leq m\leq a_1a_2\cdots a_{k-1}-1\\ d|r-ma_k}}N_{\{\frac{a_1}{d},\frac{a_2}{d},\cdots ,\frac{a_{k-1}}{d}\}}^p\left(\frac{a_1a_2\cdots a_kl}{d}+\frac{r-ma_k}{d}\right)\\
&=\sum_{i=1}^{a_1a_2\cdots a_{k-1}}p_A(a_1a_2\cdots a_kl+r-ia_k)
\end{align*}
\begin{align}\label{eqn2}
\ \ \ \ \ \ \ \ \ \  +\sum_{\substack{0\leq m\leq a_1a_2\cdots a_{k-1}-1\\ d|r-ma_k}}N_{\{\frac{a_1}{d},\frac{a_2}{d},\cdots ,\frac{a_{k-1}}{d}\}}^p\left(\frac{a_1}{d}\frac{a_2}{d}\cdots \frac{a_{k-1}}{d}\frac{d^{k-1}a_kl}{d}+\frac{r-ma_k}{d}\right).
\end{align}
By Lemma \ref{lemma2}, $p_A(a_1a_2\cdots a_kl+r-ia_k)$ is a polynomial in $l$ or $l-1$ (depending respectively on the bounds: $r-ia_k\geq 0$ or $r-ia_k<0$) of degree $k-1$ with the leading coefficient $\frac{(a_1a_2\cdots a_k)^{k-2}}{(k-1)!}$. In both the cases, $p_A(a_1a_2\cdots a_kl+r-ia_k)$ is  a polynomial in $l$ of degree $k-1$ with the leading coefficient $\frac{(a_1a_2\cdots a_k)^{k-2}}{(k-1)!}$. Consequently, $\sum_{i=1}^{a_1a_2\cdots a_{k-1}}p_A(a_1a_2\cdots a_kl+r-ia_k)$ is a polynomial in $l$ of degree $k-1$ with the leading coefficient $(a_1a_2\cdots a_{k-1})\frac{(a_1a_2\cdots a_k)^{k-2}}{(k-1)!}=\frac{(a_1a_2\cdots a_{k-1})^{k-1}a_k^{k-2}}{(k-1)!}$.

Since $\gcd(d,a_k)=1$, the equation $a_km\equiv r\pmod d$ has a unique solution modulo $d$. Consequently, there exists a unique integer $m_t\in\{td+0,td+1,\ldots,td+(d-1)\}$ such that $a_km_t\equiv r\pmod d$ for each $t\in\{0,1,\ldots , a_1a_2\cdots a_{k-1}-1\}$. In view of division algorithm, one can find a unique pair of integers $(q_t,r_t)$ such that $\frac{r-m_ta_k}{d}=\frac{a_1a_2\cdots a_{k-1}}{d^{k-1}}q_t+r_t$. Based on these observations, we can write 
\[\sum_{\substack{0\leq m\leq a_1a_2\cdots a_{k-1}-1\\ d|r-ma_k}}N_{\{\frac{a_1}{d},\frac{a_2}{d},\cdots ,\frac{a_{k-1}}{d}\}}^p\left(\frac{a_1}{d}\frac{a_2}{d}\cdots \frac{a_{k-1}}{d}\frac{d^{k-1}a_kl}{d}+\frac{r-ma_k}{d}\right)\ \ \ \ \ 
\]
\[=\sum_{t=0}^{\frac{a_1a_2\cdots a_{k-1}}{d}-1}N_{\{\frac{a_1}{d},\frac{a_2}{d},\cdots ,\frac{a_{k-1}}{d}\}}^p\left(\frac{a_1}{d}\frac{a_2}{d}\cdots \frac{a_{k-1}}{d}\left(d^{k-2}a_kl+q_t\right)+r_t\right).
\]
Since $\gcd\left(\frac{a_1}{d},\frac{a_2}{d},\cdots ,\frac{a_{k-1}}{d}\right)=1$, by induction assumption, we have that 
\[
N_{\{\frac{a_1}{d},\frac{a_2}{d},\cdots ,\frac{a_{k-1}}{d}\}}^p\left(\frac{a_1}{d}\frac{a_2}{d}\cdots \frac{a_{k-1}}{d}\left(d^{k-2}a_kl+q_t\right)+r_t\right)
\]
is a polynomial in $d^{k-2}a_kl+q_t$ of degree $k-1$ with the leading coefficient 
\[\frac{\sum_{s=1}^{k-1}\left(\frac{a_s}{d}\right)^{k-3}\prod_{\substack{j\neq s\\ 0\leq j\leq k-1}}\left(\frac{a_j}{d}\right)^{k-2}}{(k-1)!}.
\]
Consequently, $N_{\{\frac{a_1}{d},\frac{a_2}{d},\cdots ,\frac{a_{k-1}}{d}\}}^p\left(\frac{a_1}{d}\frac{a_2}{d}\cdots \frac{a_{k-1}}{d}\left(d^{k-2}a_kl+q_t\right)+r_t\right)$ is a polynomial in $l$ of degree $k-1$ with the leading coefficient \[\frac{(a_kd^{k-2})^{k-1}}{d^{k-3}d^{(k-2)(k-2)}}\frac{\sum_{s=1}^{k-1}a_s^{k-3}\prod_{\substack{j\neq s\\ 0\leq j\leq k-1}}a_j^{k-2}}{(k-1)!}.\]
This gives that \[\sum_{t=0}^{\frac{a_1a_2\cdots a_{k-1}}{d}-1}N_{\{\frac{a_1}{d},\frac{a_2}{d},\cdots ,\frac{a_{k-1}}{d}\}}^p\left(\frac{a_1}{d}\frac{a_2}{d}\cdots \frac{a_{k-1}}{d}\left(d^{k-2}a_kl+q_i\right)+r_i\right)\]
is a polynomial in $l$ of degree $k-1$ with the leading coefficient 
\begin{flushleft}
$\frac{a_1a_2\cdots a_{k-1}}{d}\frac{(a_kd^{k-2})^{k-1}}{d^{k-3}d^{(k-2)(k-2)}}\frac{\sum_{s=1}^{k-1}a_s^{k-3}\prod_{\substack{j\neq s\\ 0\leq j\leq k-1}}a_j^{k-2}}{(k-1)!}
=\frac{a_k^{k-1}\sum_{s=1}^{k-1}a_s^{k-2}\prod_{\substack{j\neq s\\ 0\leq j\leq k-1}}a_j^{k-1}}{(k-1)!}$
\end{flushleft}
\begin{flushright} 
$=\frac{\sum_{s=1}^{k-1}a_s^{k-2}\prod_{\substack{j\neq s\\ 0\leq j\leq k}}a_j^{k-1}}{(k-1)!}.\ \ \ \ \ \ \ \ \ \ \ \ $
\end{flushright}
Now in view of (\ref{eqn2}), we get $N_A^p(a_1a_2\cdots a_kl+r)-N_A^p(a_1a_2\cdots a_k(l-1)+r)$ as a polynomial in $l$ of degree $k-1$ with the leading coefficient
\begin{align*}
\frac{\sum_{s=1}^{k-1}a_s^{k-2}\prod_{\substack{j\neq s\\ 0\leq j\leq k}}a_j^{k-1}}{(k-1)!}+\frac{a_k^{k-2}(a_1a_2\cdots a_{k-1})^{k-1}}{(k-1)!}=\frac{\sum_{i=1}^{k}a_i^{k-2}\prod_{\substack{j\neq i\\ 0\leq j\leq k}}a_j^{k-1}}{(k-1)!}.
\end{align*} 
Let $c_k$ be the leading coefficient of $N_A^p(a_1a_2\cdots a_kl+r)$. Then from the above observations, we have
\[kc_k=\frac{\sum_{i=1}^{k}a_i^{k-2}\prod_{\substack{j\neq i\\ 0\leq j\leq k}}a_j^{k-1}}{(k-1)!},
\] 
which gives 
\[c_k=\frac{\sum_{i=1}^{k}a_i^{k-2}\prod_{\substack{j\neq i\\ 0\leq j\leq k}}a_j^{k-1}}{k!}.
\] 
Hence, the main assertion is true for $k$. Then the targeted estimate follows immediately while adhering to the discussion at the starting part of this proof. 
\end{proof}

\end{document}
